\documentclass[11pt,a4paper]{article}
\usepackage{fontspec}
\usepackage[margin=2.5cm]{geometry}
\usepackage{amsmath,amssymb,amsthm,mathtools}
\usepackage{booktabs,enumitem,tabularx,array}
\usepackage{float}
\usepackage{tikz}
\usetikzlibrary{positioning,arrows.meta,calc}
\usepackage[colorlinks=true,linkcolor=blue!60!black,citecolor=blue!60!black,urlcolor=blue!60!black]{hyperref}
\newfontfamily\lmfont{Latin Modern Roman}
\newcommand{\lstroke}{{\lmfont ł}}
\setmainfont{Noto Serif CJK KR}[ItalicFont={Noto Serif CJK KR},ItalicFeatures={FakeSlant=0.18},BoldFont={Noto Serif CJK KR Bold},BoldItalicFont={Noto Serif CJK KR Bold},BoldItalicFeatures={FakeSlant=0.18}]
\renewcommand{\arraystretch}{1.25}
\linespread{1.12}
\setlength{\parskip}{0.35em}
\setlength{\emergencystretch}{10em}
\tolerance=2000
\frenchspacing

\renewcommand{\abstractname}{초록}
\renewcommand{\refname}{참고문헌}
\renewcommand{\proofname}{증명}
\renewcommand{\figurename}{그림}

\renewcommand{\qedsymbol}{$\blacksquare$}

\theoremstyle{plain}
\newtheorem{theorem}{정리}[section]
\newtheorem{lemma}[theorem]{보조정리}
\newtheorem{corollary}[theorem]{따름정리}
\newtheorem{proposition}[theorem]{명제}
\theoremstyle{definition}
\newtheorem{definition}[theorem]{정의}
\newtheorem{remark}[theorem]{주석}
\newtheorem{example}[theorem]{예}

\newcommand{\R}{\mathbb R}
\newcommand{\Z}{\mathbb Z}
\newcommand{\N}{\mathbb N}
\newcommand{\AM}{A_M}
\newcommand{\kerZ}{\ker_{\Z}\!A_M}
\newcommand{\kerR}{\ker_{\R}\!A_M}
\newcommand{\sq}{L_{\square}}
\newcommand{\bs}{\mathcal B}
\newcommand{\sgn}{\operatorname{sgn}}
\newcommand{\rank}{\operatorname{rank}}
\newcommand{\cl}{\operatorname{cl}}
\newcommand{\ind}{\mathbf 1}
\newcommand{\colex}{<_{\mathrm{colex}}}
\newcommand{\LM}{Lasoń--Micha\lstroke{}ek}

\title{매트로이드의 정사각형 기저 정리}
\author{김영수}
\date{}

\begin{document}
\maketitle

\begin{abstract}
\noindent White 추측은 모든 매트로이드의 토릭 아이디얼이 대칭 교환 이항식들로 생성된다는 주장, 즉 교환 이동들이 Markov basis를 이룬다는 주장이며, 일반적으로는 성립하지 않는다. 본 논문은 그 격자 판이 계수가 $1$ 이상인 모든 매트로이드에서 성립함을 증명한다. 구체적으로, 대칭 교환 이항식의 지수벡터인 \emph{정사각형 관계} $e_{B+i}-e_{B+j}-e_{B'+i}+e_{B'+j}$들은 기저 결합 행렬 $\AM$의 정수 핵 $\kerZ$의 $\Z$\nobreakdash-기저를 항상 포함한다. 이는 다음과 동치이다. 임의의 체 위에서 토릭 아이디얼 $I_M$은 변수들을 가역으로 만들면 $\operatorname{codim} I_M$개의 대칭 교환 이항식의 완전 교차가 된다. 증명은 구성적이고 조합적이며, 컴퓨터 계산을 쓰지 않는다. 직합 분해, 원소 하나를 삭제할 때의 완전열, 그리고 매트로이드의 고정된 원소를 지나는 circuit들의 연결성을 결합하며, 마지막 성질은 Lehman 정리와 밀접하게 관련되어 있다. 균등 매트로이드에 대해서는 단위삼각 정사각형 기저를 명시적으로 제시한다. 응용으로, 두 실수 행렬이 가역 행 변환과 가역 열 스케일링으로 서로 옮겨지는지를 판정한다. 판정 조건은 정사각형 기저의 원소마다 하나씩 주어지는, Plücker 비의 부호와 절댓값에 대한 정사각형 조건이다. 부호 조건에는 이 기저가 단순한 $\R$\nobreakdash-기저가 아니라 $\Z$\nobreakdash-기저라는 사실이 필요하다.

\vspace{1.2em}\noindent\textbf{키워드.} 매트로이드, 토릭 아이디얼, 기저 교환, 격자 기저, White 추측, Plücker 좌표

\smallskip\noindent\textbf{수학 주제 분류.} 05B35, 13F65, 14M25
\end{abstract}

%====================================================================
\section{서론}\label{sec:intro}
White는 모든 매트로이드 $M$의 토릭 아이디얼 $I_M$이 대칭 기저 교환에서 오는 이차 이항식들로 생성된다고 추측하였다~\cite{W80}. 이 추측은 계수 $3$ 이하의 매트로이드~\cite{Ka10}, 그래프 매트로이드~\cite{Bl08}, 격자 경로 매트로이드~\cite{Sc11}, sparse paving 매트로이드~\cite{Bo13}, paving 매트로이드~\cite{YY25}, strongly base orderable 매트로이드~\cite{LM14}에 대해 알려져 있다. 또한 모든 매트로이드에서 대칭 교환 이항식들이 생성하는 아이디얼을 무관 아이디얼(모든 변수가 생성하는 아이디얼)로 포화하면 토릭 아이디얼과 같아지고~\cite{LM14}, 계수를 고정하면 추측이 높은 차수에서 성립한다~\cite{La21}. 이 추측은 기저 하나를 더하거나 빼도 보존되며~\cite{HMW25}, 원소가 $7$개 이하인 매트로이드 가운데 두 개를 제외한 모든 경우에 Gröbner 기저가 계산되었다~\cite{HHKS20}. regular 매트로이드와 split 매트로이드에서는 기저 쌍 사이의 교환 거리도 연구되었다~\cite{BMS24,BS24}. 2026년에 Larson이 White 추측의 반례를 제시하였고~\cite{Lar26}, Fano 또는 쌍대 Fano minor를 가진 매트로이드에는 이차 Gröbner 기저가 없음이 밝혀졌다~\cite{BCLLM26,DFMR26}.

대수 통계학의 용어를 빌리면~\cite{DS98}, White 추측은 교환 이동(exchange move)들이 $I_M$에 대응하는 격자의 \emph{Markov basis}를 이룬다는 주장이며, 이는 일반적으로 성립하지 않는다. 이와 달리 격자 수준에서는 이동들의 집합이 \emph{lattice basis}, 즉 같은 격자의 $\Z$\nobreakdash-기저를 포함할 것을 요구할 수 있다. 본 논문은 이 요구가 계수가 $1$ 이상인 모든 매트로이드에서 성립함을 증명한다(정리~\ref{thm:main}). 이 진술은 $I_M$ 자체의 생성원이 아니라 $I_M$의 격자에 관한 것이므로, 위의 모든 결과와 양립한다(주석~\ref{rem:markov}, \ref{rem:white}).

$M$을 $[n]=\{1,\dots,n\}$ 위의 계수 $k$인 매트로이드, $\bs$를 그 기저들의 집합, $N=|\bs|$라 하자. \emph{기저 결합 행렬} $\AM\in\{0,1\}^{n\times N}$의 $I$에 해당하는 열은 기저 $I$의 지시벡터 $\ind_I$이다. 그 정수 핵
\[
\kerZ=\{u\in\Z^{\bs}:\AM u=0\}
\]
은 계수 $N-\rank\AM$인 자유 아벨 군이며, $M$의 토릭 아이디얼 $I_M$, 즉 사상 $y_I\mapsto\prod_{i\in I}x_i$의 핵에 대응하는 격자이다~\cite{W77,St96}.

기저 교환은 이 격자의 자연스러운 원소를 만든다. $i\ne j$이고, $B,B'$은 $i,j\notin B\cup B'$인 서로 다른 $(k-1)$\nobreakdash-집합이라 하자. $B+i$, $B+j$, $B'+i$, $B'+j$가 모두 기저이면 $\Z^n$에서 $\ind_{B+i}-\ind_{B+j}=e_i-e_j=\ind_{B'+i}-\ind_{B'+j}$이다. 따라서 $I\in\bs$에 해당하는 $\Z^{\bs}$의 표준 기저 벡터를 $e_I$로 쓰면
\[
z_{i,j;B,B'}=e_{B+i}-e_{B+j}-e_{B'+i}+e_{B'+j}\in\kerZ
\]
이다. 이 $z_{i,j;B,B'}$을 \emph{정사각형 관계}라 한다. 정사각형 관계는 White 추측~\cite{W80}에 나오는 대칭 교환 이항식 $y_{B_1}y_{B_2}-y_{B_1-a+b}\,y_{B_2-b+a}$의 지수벡터와 정확히 일치한다(보조정리~\ref{lem:sym}).

\begin{theorem}[정사각형 기저 정리]\label{thm:main}
계수가 $1$ 이상인 모든 매트로이드 $M$에 대해, 격자 $\kerZ$은 $N-\rank\AM$개의 정사각형 관계로 이루어진 $\Z$\nobreakdash-기저를 가진다.
\end{theorem}

계수 조건은 필요하다. $\rank M=0$이면 $\kerZ=\Z$이지만 정사각형은 하나도 없다. 그림~\ref{fig:u24}은 자명하지 않은 균등 매트로이드 가운데 가장 작은 예를 보여 준다.

\begin{figure}[H]
\centering
\newcommand{\octa}[3]{%
\begin{tikzpicture}[scale=0.95,v/.style={circle,fill=black,inner sep=1.1pt},lab/.style={font=\footnotesize}]
\coordinate (T) at (0,1.75);  \coordinate (D) at (0,-1.75);
\coordinate (L) at (-1.45,0); \coordinate (R) at (1.45,0);
\coordinate (K) at (0.55,0.42); \coordinate (F) at (-0.55,-0.42);
#1
\draw[gray!60,dashed] (T)--(K) (D)--(K) (L)--(K) (R)--(K);
\draw (T)--(L) (T)--(R) (T)--(F) (D)--(L) (D)--(R) (D)--(F) (L)--(F) (F)--(R);
#2
\foreach \p in {T,D,L,R,K,F} \node[v] at (\p) {};
\node[lab,above] at (T) {$12$}; \node[lab,below] at (D) {$34$};
\node[lab,left] at (L) {$24$}; \node[lab,right] at (R) {$13$};
\node[lab,above right=-1pt] at (K) {$14$}; \node[lab,below left=-1pt] at (F) {$23$};
\node[font=\small,align=center] at (0,-2.55) {#3};
\end{tikzpicture}}
\octa{\fill[blue!20] (L)--(K)--(R)--(F)--cycle;}{\draw[blue!70!black,very thick] (L)--(F)--(R); \draw[blue!70!black,very thick,dashed] (R)--(K)--(L);}{$z_1=e_{13}-e_{14}+e_{24}-e_{23}$}\hfill
\octa{\fill[red!18] (T)--(K)--(D)--(F)--cycle;}{\draw[red!70!black,very thick] (T)--(F)--(D); \draw[red!70!black,very thick,dashed] (D)--(K)--(T);}{$z_2=e_{12}-e_{14}+e_{34}-e_{23}$}\hfill
\octa{\fill[green!20] (T)--(R)--(D)--(L)--cycle;}{\draw[green!40!black,very thick] (T)--(R)--(D)--(L)--cycle;}{$z_3=e_{12}-e_{13}+e_{34}-e_{24}$}
\caption{$U_{2,4}$의 기저 다면체는 정팔면체이며, 꼭짓점은 $|I|=2$인 여섯 기저 $I\subset[4]$에 대응한다. 마주 보는 꼭짓점은 서로 여집합이다. 부호를 무시하면 세 개의 적도 정사각형이 정사각형 관계의 전부이다(부호는 정사각형을 따라 번갈아 바뀐다). 여기서 $\rank\kerZ=6-4=2$이고 $z_1=z_2-z_3$이므로, 세 정사각형 중 어느 두 개든 $\kerZ$의 $\Z$\nobreakdash-기저를 이룬다.}
\label{fig:u24}
\end{figure}

\medskip\noindent\textbf{기저가 왜 중요한가.} 정리~\ref{thm:main}은 계수가 $1$ 이상인 매트로이드에 대해 이미 알려진 두 사실을 강화한다. $\R$ 위에서는 정사각형들이 $\kerR$을 생성하며, 이는 아래 정리~\ref{thm:sat}에서도 따라 나온다. 따라서 정사각형 가운데 $N-\rank\AM$개를 골라 $\R$\nobreakdash-기저를 만들 수 있다. 그러나 격자의 $\R$\nobreakdash-기저가 그 격자를 생성한다는 보장은 없다. $\Z$ 위에서는 정사각형들이 $\kerZ$을 생성하며, 이는 \LM{} 정리~\cite{LM14}에서 따라 나온다(명제~\ref{prop:lattice}). 그러나 자유 아벨 군의 생성 집합이 기저를 포함한다는 보장도 없다. 예를 들어 $\{2,3\}$은 $\Z$를 생성하지만, 그 어느 부분집합도 $\Z$의 기저가 아니다. 정리~\ref{thm:main}은 두 성질을 동시에 준다. 즉 $\kerZ$을 $\Z$ 위에서 생성하면서 크기가 가능한 최소인 정사각형 집합이 존재한다. 이로부터 다음이 따라 나온다.

\begin{enumerate}[label=(\roman*),leftmargin=2em,itemsep=2pt]
\item \emph{토릭 아이디얼.} $c=N-\rank\AM$이라 하자. 이는 $I_M$의 여차원이다. 정리~\ref{thm:main}은 임의의 체 $K$ 위에서
\[
I_M\,K[y^{\pm1}]=(f_1,\dots,f_c)\,K[y^{\pm1}],\qquad I_M=(f_1,\dots,f_c):\Big(\prod_{I\in\bs}y_I\Big)^{\infty}
\]
을 만족하는 $c$개의 대칭 교환 이항식 $f_1,\dots,f_c$가 존재한다는 것과 동치이다(따름정리~\ref{cor:toric}). Krull 높이 정리에 의해 생성원이 $c$개보다 적으면 두 등식 중 어느 것도 성립할 수 없다. 기하적으로 말하면, $M$의 토릭 다양체 중 토러스 $(K^\times)^{\bs}$ 안에 놓인 부분은 $c$개의 대칭 교환 이차식의 완전 교차이다. 반면 $I_M$ 자체가 완전 교차가 되는 경우는 드물다. $M$에 루프와 코루프가 없고 $c\ge1$이면, 이는 원소가 $4$개인 세 매트로이드에서만 일어난다~\cite{GR15}.
\item \emph{국소 증명서.} $\Theta$를 임의의 아벨 군, 예를 들어 $\Z/2$나 $\R$이라 하자. 함수 $\chi:\bs\to\Theta$가 정사각형 기저를 이루는 $c$개의 정사각형에 대해 정사각형 조건 $\chi(B+i)-\chi(B+j)-\chi(B'+i)+\chi(B'+j)=0$을 만족하면, 모든 $u\in\kerZ$에 대해 $\sum_I u_I\chi(I)=0$이다. 사상 $u\mapsto\sum_Iu_I\chi(I)$는 생성 집합에서 $0$이 되는 준동형사상이기 때문이다. $\Z/2$ 같은 꼬임 군(torsion group)에서는 정사각형으로 된 $\R$\nobreakdash-기저만으로는 일반적으로 충분하지 않은데, 그러한 기저가 유한 지수의 진부분격자만 생성할 수 있기 때문이다(주석~\ref{rem:index}). 또한 기저가 아닌 정사각형 생성 집합은 $c$개보다 많은 원소를 가진다. $\Theta=\Z/2$인 경우와 $\Theta=\R$인 경우를 함께 쓰면, 두 실수 행렬이 가역 행 변환과 가역 열 스케일링으로 서로 옮겨지는지에 대한 판정법을 얻는다(\S\ref{sec:criterion}).
\end{enumerate}


\medskip\noindent\textbf{증명의 개요.} 증명의 흐름은 그림~\ref{fig:map}에 정리되어 있으며, \LM{} 정리는 사용하지 않는다. 먼저 직합을 이용하여 문제를 연결 매트로이드의 경우로 환원한다(\S\ref{sec:dsum}). 다음으로 연결 매트로이드 $M$의 원소 $v$에 대해 $M\setminus v$가 연결이라 하자. 이때 완전열 $0\to\ker_\Z A_{M\setminus v}\to\kerZ\to\Gamma\to0$을 이용하면, $[n]-v$ 위의 어떤 그래프가 연결이면 $M\setminus v$의 정사각형 기저를 $M$으로 들어 올릴 수 있다(\S\ref{sec:conn}). 이 그래프가 연결일 필요충분조건은 $v$를 지나는 $M$의 cocircuit들이 나머지 모든 원소를 이어 주는 것이다(명제~\ref{prop:phi}). 이러한 cocircuit은 바로 $v$를 지나는 $M^*$의 circuit이다. \S\ref{sec:phiproof}에서는 연결 매트로이드 $P$에서 $P/v$가 연결이면 $v$를 지나는 circuit들이 나머지 모든 원소를 이어 준다는 것을 증명한다(정리~\ref{thm:phi}). $P=M^*$로 두면 $P/v=(M\setminus v)^*$이므로, 이 가정은 정확히 $M\setminus v$가 연결이라는 조건이 된다. 끝으로 Tutte 정리~\cite[\S4.3]{Ox11}에 따르면 $M\setminus v$와 $M/v$ 중 적어도 하나는 연결이다. $M/v$가 연결이면 $M^*\setminus v=(M/v)^*$이 연결이므로, 쌍대성(보조정리~\ref{lem:dual})에 의해 $M^*$에 앞의 논증을 적용할 수 있다. 이로써 귀납이 완성된다. 균등 매트로이드에 대해서는 \S\ref{sec:uniform}에서 명시적인 기저를 제시한다.

\begin{figure}[tbp]
\centering
\begin{tikzpicture}[
  box/.style={draw,rounded corners=3pt,align=center,font=\small,inner sep=5pt,text width=5.4cm,minimum height=1.35cm},
  key/.style={box,fill=blue!8,draw=blue!60!black,thick},
  arr/.style={-{Stealth[length=2.2mm]},thick}]
\node[key] (phi) {$P$와 $P/v$가 연결이면\\ $v$를 지나는 circuit들이 $E(P)-v$를 잇는다\\ \footnotesize 정리~\ref{thm:phi}, \S\ref{sec:phiproof}};
\node[box,right=1.0cm of phi] (H) {$M/v$의 기저들 위의\\ 그래프 $H$가 연결이다\\ \footnotesize 명제~\ref{prop:phi} (쌍대성)};
\node[box,below=0.9cm of H] (lift) {완전열과 spanning tree:\\ $M\setminus v$의 정사각형 기저를\\ $M$으로 들어 올림\\ \footnotesize 정리~\ref{thm:exact}, \ref{thm:tree}};
\node[box,below=0.9cm of phi] (tutte) {Tutte: $M\setminus v$ 또는 $M/v$가 연결;\\ 정사각형 기저의 존재는 $M$과 $M^*$에서 동치\\ \footnotesize Tutte 정리, 보조정리~\ref{lem:dual}};
\node[box,below=0.9cm of lift] (conn) {계수가 $1$ 이상인\\ 모든 연결 매트로이드는\\ 정사각형 기저를 가진다\\ \footnotesize 정리~\ref{thm:cond}};
\node[box,below=0.9cm of tutte] (dsum) {직합은 정사각형 기저를\\ 보존한다\\ \footnotesize 명제~\ref{prop:dsum}};
\node[key,minimum height=0.9cm] (main) at ($(dsum.south)!0.5!(conn.south)+(0,-1.15)$) {정사각형 기저 정리~\ref{thm:main}};
\draw[arr] (phi)--(H);
\draw[arr] (H)--(lift);
\draw[arr] (lift)--(conn);
\draw[arr] (tutte)--(conn);
\draw[arr] (conn.south) |- (main.east);
\draw[arr] (dsum.south) |- (main.west);
\end{tikzpicture}
\caption{증명의 구조. 왼쪽 위의 음영 상자는 증명의 중심이 되는 매트로이드 논증이다. \S\ref{sec:phiproof}에서 자기완결적인 증명을 주고, 주석~\ref{rem:lehman}에서 Lehman 정리로부터 유도한다.}
\label{fig:map}
\end{figure}

\medskip\noindent\textbf{구성.} \S\ref{sec:prelim}에서 기호를 정한다. \S\ref{sec:dsum}--\S\ref{sec:phiproof}에서 정리~\ref{thm:main}을 증명한다. \S\ref{sec:uniform}에서는 균등 매트로이드를 다룬다. \S\ref{sec:apps}에서는 토릭 아이디얼에 대한 귀결과 가역 열 스케일링에 대한 응용을 제시하고, \S\ref{sec:caveats}에서는 \LM{} 정리와의 관계와 열린 문제를 논의한다.

%====================================================================
\section{준비}\label{sec:prelim}
매트로이드 용어는 Oxley의 책~\cite{Ox11}을 따른다. 집합 $I$와 원소 $a$에 대해 $I+a=I\cup\{a\}$, $I-a=I\setminus\{a\}$로 쓴다. 모든 기저의 크기가 $k$이므로 $\AM^{T}\ind_n=k\,\ind_{\bs}$이고, 따라서 $k\ge1$이면 모든 $u\in\kerZ$의 좌표합은 $0$이다. 체 $K$ 위에서 $K[y]=K[y_I:I\in\bs]$라 하면, 토릭 아이디얼 $I_M\subseteq K[y]$는 사상 $K[y]\to K[x_1,\dots,x_n]$, $y_I\mapsto\prod_{i\in I}x_i$의 핵이며
\[
I_M=\big(y^{u^+}-y^{u^-}:u\in\kerZ\big)
\]
이다~\cite{St96}. 여기서 $u=u^+-u^-$는 $u$를 지지집합이 서로소인 두 벡터 $u^+,u^-\in\N^{\bs}$의 차로 쓴 것이고, $y^w=\prod_Iy_I^{w_I}$이다. $K[y]/I_M$의 Krull 차원은 $\rank\AM$이므로 $I_M$의 높이는 $N-\rank\AM$이다.

\begin{definition}\label{def:square}
서로 다른 $i,j\in[n]$과, $i$와 $j$를 포함하지 않는 서로 다른 $(k-1)$\nobreakdash-부분집합 $B,B'$을 생각하자. $B+i,B+j,B'+i,B'+j$가 모두 기저이면 벡터 $z_{i,j;B,B'}=e_{B+i}-e_{B+j}-e_{B'+i}+e_{B'+j}\in\Z^{\bs}$를 \emph{정사각형 관계} 또는 간단히 \emph{정사각형}이라 한다. 정사각형들이 생성하는 $\Z^{\bs}$의 부분격자를 $\sq(M)$이라 쓰고, $\kerZ$이 정사각형들로 이루어진 $\Z$\nobreakdash-기저를 가지면 $M$이 \emph{정사각형 기저를 가진다}고 한다.
\end{definition}
네 지시벡터가 $\ind_{B+i}-\ind_{B+j}-\ind_{B'+i}+\ind_{B'+j}=0$을 만족하므로 항상 $\sq(M)\subseteq\kerZ$이다. \emph{대칭 교환}은 기저 $B_1,B_2$와 원소 $a\in B_1\setminus B_2$, $b\in B_2\setminus B_1$로서 $B_1-a+b$와 $B_2-b+a$가 모두 기저인 것을 말하며, 그 이항식은 $y_{B_1}y_{B_2}-y_{B_1-a+b}\,y_{B_2-b+a}$이다. 대칭 교환 정리~\cite{Br69}에 의해, 서로 다른 기저 $B_1,B_2$와 임의의 $a\in B_1\setminus B_2$에 대해 이러한 $b$가 항상 존재한다.

\begin{lemma}\label{lem:sym}
정사각형은 $0$이 아닌 대칭 교환 이항식의 지수벡터와 정확히 일치한다.
\end{lemma}
\begin{proof}
정사각형이 주어지면 $B_1=B+i$, $B_2=B'+j$, $a=i$, $b=j$로 둔다. 그러면 $B_1-a+b=B+j$, $B_2-b+a=B'+i$이고, $i,j\notin B\cup B'$은 $a\in B_1\setminus B_2$, $b\in B_2\setminus B_1$과 동치이다. 역으로 대칭 교환이 주어지면 $B=B_1-a$, $B'=B_2-b$, $i=a$, $j=b$로 둔다. 이항식이 $0$이 되는 것은 정확히 $B=B'$일 때이다.
\end{proof}

\begin{proposition}[경계 경우]\label{prop:edge}
\begin{enumerate}[label=(\alph*),itemsep=1pt]
\item $\rank M=0$이면 $\bs=\{\varnothing\}$이고 $\kerZ=\Z$이지만 정사각형은 없다.
\item 루프를 삭제해도 $\kerZ$과 정사각형들의 집합은 바뀌지 않는다.
\item $M=M'\oplus\{c\}$이고 $c$가 코루프이며 $\rank M'\ge1$이면, $I'\mapsto I'+c$는 정사각형을 정사각형으로 보내는 동형사상 $\ker_\Z A_{M'}\cong\kerZ$을 유도한다. $\rank M'=0$이면 $\kerZ=0$이다.
\end{enumerate}
\end{proposition}
\begin{proof}
(a)와 (b)는 자명하다. (c) $c$에 해당하는 $\AM$의 행은 모든 성분이 $1$인 벡터이고, 이 행이 주는 조건 $\sum f=0$은 $\rank M'\ge1$일 때 $M'$의 행들에서 이미 따라 나온다. 남은 것은 $M$의 모든 정사각형 $z_{i,j;B,B'}$에서 $c\in B\cap B'$임을 보이는 것이다. $c\notin B\cup B'$이면 $B+i$와 $B+j$가 모두 코루프 $c$를 포함해야 하므로 $i=c=j$가 되어 $i\ne j$에 모순이다. $c\in B\setminus B'$이면 $B'+i$가 코루프 $c$를 포함해야 하므로 $i=c\in B$가 되어 모순이다. $c\in B'\setminus B$인 경우도 대칭적이다.
\end{proof}

\begin{lemma}[쌍대성]\label{lem:dual}
$1\le k\le n-1$이면 $e_I\mapsto e_{[n]\setminus I}$는 동형사상 $\kerZ\to\ker_\Z A_{M^*}$을 유도하며, 이 사상은 정사각형 $z_{i,j;B,B'}$을 $-z_{i,j;C,C'}$으로 보낸다. 여기서 $C=[n]\setminus(B+i+j)$, $C'=[n]\setminus(B'+i+j)$이다. 따라서 $M$이 정사각형 기저를 가질 필요충분조건은 $M^*$이 정사각형 기저를 가지는 것이다.
\end{lemma}
\begin{proof}
$f\in\kerZ$이면 $\sum f=0$이고, $A_{M^*}$의 $i$번째 행이 주는 조건은 $\sum_{I\not\ni i}f(I)=\sum f-(\AM f)_i=0$이다. 역방향은 $M^{**}=M$을 써서 $M^*$에 같은 논증을 적용하면 된다. 정사각형에 대해서는 $[n]\setminus(B+i)=C+j$, $[n]\setminus(B+j)=C+i$이고, $B'$에 대해서도 마찬가지이다.
\end{proof}

%====================================================================
\section{직합}\label{sec:dsum}
바탕집합이 $E_1,E_2$인 $M=M_1\oplus M_2$를 생각하자. $M$의 기저는 $I_t\in\bs_t=\bs(M_t)$에 대한 $I_1\sqcup I_2$들이므로 $V_t=\Z^{\bs_t}$에 대해 $\Z^{\bs}=V_1\otimes V_2$이다. $N_t=|\bs_t|$, $\rho_t=\rank A_{M_t}$, $K_t=\ker_\Z A_{M_t}$, $k_t=\rank M_t$로 쓰고, $V_t^0\subseteq V_t$를 좌표합이 $0$인 벡터들의 집합이라 하자.

\begin{proposition}\label{prop:dsum}
$k_1,k_2\ge1$이고 $M_1,M_2$가 정사각형 기저 $\mathcal R_1,\mathcal R_2$를 가지면, $M_1\oplus M_2$도 정사각형 기저를 가진다.
\end{proposition}
\begin{proof}
\emph{1단계(핵).} $f\in V_1\otimes V_2$에 대해 $\operatorname{row}(f)(I_1)=\sum_{I_2}f(I_1,I_2)$, $\operatorname{col}(f)(I_2)=\sum_{I_1}f(I_1,I_2)$로 둔다. $E_1$에 해당하는 $\AM$의 행들은 $A_{M_1}\operatorname{row}(f)$를, $E_2$에 해당하는 행들은 $A_{M_2}\operatorname{col}(f)$를 주므로
\[
\kerZ=\{f:\operatorname{row}(f)\in K_1,\ \operatorname{col}(f)\in K_2\}
\]
이고, $k_t\ge1$이므로 $K_t\subseteq V_t^0$이다.

\emph{2단계(분해).} $I_1^0\in\bs_1$과 $I_2^0\in\bs_2$를 고정하면
\[
\kerZ=K_1\otimes e_{I_2^0}\ \oplus\ e_{I_1^0}\otimes K_2\ \oplus\ (V_1^0\otimes V_2^0)
\]
이다. 앞의 두 성분은 두 인자의 핵을 담고, 세 번째 성분은 상호작용 항들로 이루어진다. 1단계에 의해 각 성분은 $\kerZ$에 속한다. 여기서 $V_1^0\otimes V_2^0$은 정확히 $\operatorname{row}(f)=0$, $\operatorname{col}(f)=0$인 $f$들의 집합이다. $f\in\kerZ$에 대해 $f'=f-\operatorname{row}(f)\otimes e_{I_2^0}-e_{I_1^0}\otimes\operatorname{col}(f)$로 두면, $\sum\operatorname{row}(f)=\sum\operatorname{col}(f)=0$이므로 $\operatorname{row}(f')=\operatorname{col}(f')=0$이다. $\operatorname{row}$와 $\operatorname{col}$을 적용하면 합이 직합임을 알 수 있다.

\emph{3단계($V_1^0/K_1$의 기저).} 격자 $\Lambda_1=A_{M_1}(V_1^0)$은 $J,J'\in\bs_1$에 대한 벡터 $\ind_J-\ind_{J'}$들로 생성된다. 기저 교환 공리에 의해 $M_1$의 임의의 두 기저는 이웃한 것끼리 원소 하나만 다른 기저들의 열로 이어지므로 인접한 쌍만으로 충분하고, $J=D+a$, $J'=D+b$이면 $e_a-e_b$를 얻는다. 어떤 $D$에 대해 $D+a,D+b\in\bs_1$일 때 $a\sim b$로 정한 $E_1$ 위의 그래프를 $G_1$이라 하고, $F_1$을 $G_1$의 spanning forest라 하자. 벡터 $e_a-e_b$ ($\{a,b\}\in F_1$)들은 $\Lambda_1$의 $\Z$\nobreakdash-기저를 이룬다. 임의의 $J_0\in\bs_1$에 대해 $A_{M_1}(V_1)=\Lambda_1+\Z\ind_{J_0}$이고 $\ind_{J_0}$의 좌표합은 $k_1\neq0$이므로 $|F_1|=\rank\Lambda_1=\rho_1-1$이다. 각 변 $\{a,b\}\in F_1$에 대해 $D_{ab}+a,D_{ab}+b\in\bs_1$인 $D_{ab}$를 고르고 $u_{ab}=e_{D_{ab}+a}-e_{D_{ab}+b}\in V_1^0$으로 둔다. 그러면 $A_{M_1}u_{ab}=e_a-e_b$이고, $A_{M_1}$이 동형사상 $V_1^0/K_1\cong\Lambda_1$을 유도하므로 잉여류(coset) $\overline{u_{ab}}$들은 $V_1^0/K_1$의 $\Z$\nobreakdash-기저를 이룬다.

\emph{4단계(기저).} $\kerZ$의 부분군 $K_1\otimes V_2=K_1\otimes e_{I_2^0}\oplus K_1\otimes V_2^0$에 대한 몫
\[
\kerZ/(K_1\otimes V_2)\cong e_{I_1^0}\otimes K_2\ \oplus\ (V_1^0/K_1)\otimes V_2^0
\]
은 자유 군이다. $V_2^0$의 기저 $\{e_I-e_{I_2^0}\}_{I\ne I_2^0}$을 쓰면 $\kerZ$의 다음 $\Z$\nobreakdash-기저를 얻는다.
\begin{enumerate}[label=(\arabic*),itemsep=1pt]
\item $w\in\mathcal R_1$, $I\in\bs_2$에 대한 $w\otimes e_I$;
\item $s\in\mathcal R_2$에 대한 $e_{I_1^0}\otimes s$;
\item $\{a,b\}\in F_1$, $I\in\bs_2\setminus\{I_2^0\}$에 대한 $u_{ab}\otimes(e_I-e_{I_2^0})$.
\end{enumerate}
그 개수는 $N_1N_2-\rho_1-\rho_2+1=N_1N_2-\rank\AM$이다. $\AM f=(A_{M_1}\operatorname{row}(f),\,A_{M_2}\operatorname{col}(f))$이고, $f$가 $\Z^{\bs}$ 전체를 움직일 때 쌍 $(\operatorname{row}(f),\operatorname{col}(f))$는 정확히 좌표합이 같은 쌍 전체이기 때문이다.

\emph{5단계(모두 정사각형).} $w=z_{i,j;D,D'}$이면 $w\otimes e_I=z_{i,j;D\sqcup I,\,D'\sqcup I}$이고, (2)도 마찬가지이다. (3)에 대해서는 $B=D_{ab}\sqcup I$, $B'=D_{ab}\sqcup I_2^0$, $i=a$, $j=b$로 둔다. 네 집합 $(D_{ab}+a)\sqcup I$, $(D_{ab}+b)\sqcup I$, $(D_{ab}+a)\sqcup I_2^0$, $(D_{ab}+b)\sqcup I_2^0$은 기저이고, $a,b\notin B\cup B'$, $B\ne B'$이며, 해당 정사각형은 $u_{ab}\otimes(e_I-e_{I_2^0})$이다.
\end{proof}

%====================================================================
\section{연결 매트로이드}\label{sec:conn}
이 절에서 $M$은 $n\ge2$인 연결 매트로이드이다. 따라서 루프도 코루프도 없고 $1\le k\le n-1$이다. $v\in[n]$을 고정하고 $X=[n]-v$, $M_{\mathrm d}=M\setminus v$, $M_{\mathrm c}=M/v$로 둔다. $M_{\mathrm d}$의 기저는 $v$를 포함하지 않는 $M$의 기저이고, $M_{\mathrm c}$의 기저는 $J+v\in\bs$인 집합 $J\subseteq X$이다. 매트로이드 $M_0$에 대해 $\Lambda(M_0)$을 그 기저들의 지시벡터가 생성하는 격자라 하자.

\begin{lemma}\label{lem:lat}
$M_0$이 $Y$ 위의 계수 $k\ge1$인 연결 매트로이드이면 $\Lambda(M_0)=\{x\in\Z^Y:\sum x\in k\Z\}$이다. 특히 $\rank A_{M_0}=|Y|$이다.
\end{lemma}
\begin{proof}
$\subseteq$ 방향의 포함관계는 분명하다. 서로 다른 $i,j\in Y$에 대해 $i,j$를 모두 포함하는 circuit $C$가 있다. $C-j$를 기저 $I$로 확장하면 $C$는 $I$에 대한 $j$의 기본 circuit이고 $i$를 포함하므로, $I-i+j$는 기저이고 $e_j-e_i\in\Lambda(M_0)$이다. $\sum x=kq$이면 임의의 기저 $I_0$에 대해 $x-q\ind_{I_0}$의 좌표합이 $0$이므로, 이는 벡터 $e_i-e_j$들의 정수 결합이다($|Y|=1$이면 $0$이다). 마지막 주장은 우변이 $k\Z^Y$를 포함하므로 성립한다.
\end{proof}

\begin{theorem}[완전열]\label{thm:exact}
$M$과 $M_{\mathrm d}$가 연결이면 열
\[
0\to\ker_\Z A_{M_{\mathrm d}}\to\kerZ\xrightarrow{\ \pi\ }\Gamma\to0,\qquad \Gamma=\Big\{g\in\Z^{\bs(M_{\mathrm c})}:\textstyle\sum g=0\Big\}
\]
은 완전하다. 여기서 $\pi$는 $f$를 $v$를 포함하는 기저들로 제한하는 사상이다($J+v\leftrightarrow J$로 $\bs(M_{\mathrm c})$와 동일시한다).
\end{theorem}
\begin{proof}
$f=(h,g)$로 쓰자. $h$는 $v$를 포함하지 않는 기저들 위의 값, $g$는 $v$를 포함하는 기저들 위의 값이다. 그러면 $f\in\kerZ$일 필요충분조건은 $A_{M_{\mathrm d}}h+A_{M_{\mathrm c}}g=0$($X$의 행들)이고 $\sum g=0$($v$의 행)인 것이다. 따라서 $\ker\pi=\ker_\Z A_{M_{\mathrm d}}$이고 $\operatorname{im}\pi\subseteq\Gamma$이다. 역으로 $\sum g=0$이면 $A_{M_{\mathrm c}}g$의 좌표합은 $(k-1)\sum g=0$이므로 보조정리~\ref{lem:lat}에 의해 $A_{M_{\mathrm c}}g\in\Lambda(M_{\mathrm d})$이고, $A_{M_{\mathrm d}}h=-A_{M_{\mathrm c}}g$인 정수 벡터 $h$가 존재한다.
\end{proof}

서로 다른 $J,J'$이 어떤 $j\in X\setminus(J\cup J')$에 대해 $J+j$와 $J'+j$가 모두 $M_{\mathrm d}$의 기저일 때 인접하도록 $\bs(M_{\mathrm c})$ 위의 그래프 $H=H(M,v)$를 정의한다.

\begin{theorem}[spanning tree를 따른 들어 올림]\label{thm:tree}
$M$과 $M_{\mathrm d}$가 연결이고, $H$가 연결이며, $M_{\mathrm d}$가 정사각형 기저 $\mathcal R_1$을 가진다고 하자. $H$의 spanning tree $\mathcal T$의 각 변 $\{J,J'\}$에 대해 $J$와 $J'$의 순서를 하나 정해 두고 위와 같은 $j$를 하나 골라 $z_{J,J'}=z_{v,j;J,J'}=e_{J+v}-e_{J+j}-e_{J'+v}+e_{J'+j}$로 두자. 그러면 $\mathcal R_1\cup\{z_{J,J'}:\{J,J'\}\in\mathcal T\}$는 $\kerZ$의 정사각형 기저이다.
\end{theorem}
\begin{proof}
각 $z_{J,J'}$은 $M$의 정사각형이고 $\pi(z_{J,J'})=e_J-e_{J'}$이다. 연결 그래프의 spanning tree의 변 벡터들은 꼭짓점 위의 좌표합 $0$ 격자의 $\Z$\nobreakdash-기저를 이룬다(잎을 하나씩 제거하면 된다). 따라서 $\mathcal T$의 변 $e=\{J,J'\}$에 대해 $z_e=z_{J,J'}$으로 쓰면, 각 $f\in\kerZ$에 대해 $\pi(f)=\sum_e\kappa_e\pi(z_e)$인 정수 $\kappa_e$가 유일하게 존재하고, $f-\sum_e\kappa_ez_e\in\ker\pi=\langle\mathcal R_1\rangle$이다. 일차독립성은 $\pi$를 적용한 다음 $\mathcal R_1$의 독립성을 쓰면 된다. $M_{\mathrm d}$의 정사각형은 $M$의 정사각형이다.
\end{proof}

\begin{proposition}\label{prop:phi}
$M$과 $M_{\mathrm d}$가 연결이라 하고, $v$, $j$, $j'$을 모두 포함하는 $M$의 cocircuit이 있을 때 $j\sim j'$으로 정한 $X$ 위의 그래프를 $\Phi$라 하자. 그러면 $H$가 연결일 필요충분조건은 $\Phi$가 연결인 것이다.
\end{proposition}
\begin{proof}
$J\in\bs(M_{\mathrm c})$와 $j\in X\setminus J$에 대해 $J+j$가 기저일 필요충분조건은 $j\notin\cl(J)$이다. $J+v$가 기저이므로 $\cl(J)$는 $v$를 포함하지 않는 초평면이며, $v$를 포함하지 않는 모든 초평면이 이렇게 나타난다. $\mathcal N(J)=X\setminus\cl(J)$로 두자. 이는 cocircuit $[n]\setminus\cl(J)$에서 $v$를 뺀 것이다. $X$가 $M$을 생성하므로 $\mathcal N(J)$는 공집합이 아니고, 연결 매트로이드 $M$에서 $v$와 $j$는 공통의 cocircuit에 속하므로 모든 $j\in X$는 어떤 $\mathcal N(J)$에 속한다. 또한 $H$에서 $J\sim J'$일 필요충분조건은 $\mathcal N(J)\cap\mathcal N(J')\ne\varnothing$이고, $\Phi$에서 $j\sim j'$일 필요충분조건은 어떤 $J$에 대해 $j,j'\in\mathcal N(J)$인 것이다. $j\in\mathcal N(J)$일 때 $J$와 $j$를 잇는 $\bs(M_{\mathrm c})\sqcup X$ 위의 이분 그래프를 $\mathcal G_0$이라 하자. 이 그래프에는 고립점이 없다. $\mathcal G_0$의 경로는 두 쪽을 번갈아 지나므로, $\bs(M_{\mathrm c})$의 두 꼭짓점이 $\mathcal G_0$에서 이어지는 것은 $H$에서 이어지는 것과 같다. 따라서 $\mathcal G_0$이 연결일 필요충분조건은 $H$가 연결인 것이고, 같은 이유로 이는 $\Phi$가 연결인 것과도 동치이다.
\end{proof}

\begin{theorem}[고정된 원소를 지나는 circuit]\label{thm:phi}
$P$를 원소가 두 개 이상인 연결 매트로이드, $v\in E(P)$라 하자. $P/v$가 연결이면, $v$, $x$, $y$를 모두 포함하는 $P$의 circuit이 있을 때 $x\sim y$로 정한 $E(P)-v$ 위의 그래프는 연결이다. 동치로, $M$과 $M\setminus v$가 연결이면 $\Phi$는 연결이다.
\end{theorem}
두 형태가 동치인 것은 $P=M^*$로 두면 나온다. $M$의 cocircuit은 $P$의 circuit이고, $M\setminus v$가 연결일 필요충분조건은 $P/v$가 연결인 것이다. 증명은 \S\ref{sec:phiproof}에 있다.

\begin{theorem}\label{thm:cond}
계수가 $1$ 이상인 모든 연결 매트로이드는 정사각형 기저를 가진다.
\end{theorem}
\begin{proof}
$n$에 대한 귀납법을 쓴다. $n=1$이면 매트로이드는 코루프이고 $\kerZ=0$이다. $n\ge2$이고 $v\in[n]$이라 하자. Tutte 정리~\cite[\S4.3]{Ox11}에 의해 $M\setminus v$ 또는 $M/v$가 연결이다. $M\setminus v$가 연결이면 귀납 가정에 의해 정사각형 기저를 가지고, 정리~\ref{thm:phi}에 의해 $\Phi$가 연결이고 명제~\ref{prop:phi}에 의해 $H$가 연결이므로 정리~\ref{thm:tree}을 적용할 수 있다($k=1$이면 $\bs(M_{\mathrm c})=\{\varnothing\}$이고 $\Gamma=0$이다). $M/v$가 연결이면 $M^*$과 $M^*\setminus v=(M/v)^*$이 연결이므로, 첫 번째 경우에 의해 $M^*$이 정사각형 기저를 가지고, 보조정리~\ref{lem:dual}에 의해 $M$도 그러하다.
\end{proof}

\begin{proof}[정리~\ref{thm:main}의 증명]
명제~\ref{prop:edge}(b)에 의해 루프를 삭제해도 된다. 그 결과는 계수가 $1$ 이상인 연결 성분들의 직합이다. 각 성분은 정리~\ref{thm:cond}에 의해 정사각형 기저를 가지며(코루프 성분은 $\kerZ=0$), 명제~\ref{prop:dsum}로 이들을 합친다.
\end{proof}

%====================================================================
\section{고정된 원소를 지나는 circuit: 정리~\ref{thm:phi}의 증명}\label{sec:phiproof}
이 절 전체에서 $P$는 원소가 두 개 이상인 연결 매트로이드, $v\in E(P)$, $X=E(P)-v$이고, $r$과 $\cl$은 $P$의 계수 함수와 폐포 연산자이다. $P$가 연결이므로 $v$는 루프도 코루프도 아니며 $r(X)=r(P)$이다.

\medskip\noindent\textbf{설정 (S).} $X=A\sqcup B$이고 $A,B\ne\varnothing$이며, $v$를 지나는 $P$의 모든 circuit $C$는 $C-v\subseteq A$ 또는 $C-v\subseteq B$를 만족한다. $\lambda=r(A)+r(B)-r(X)$로 둔다.

정리~\ref{thm:phi}의 그래프가 비연결이면, 그 연결 성분 하나를 $A$, 나머지를 $B$로 두어 설정 (S)를 얻는다. 아래에서는 (S)에서 $\lambda=1$임을 보이고, 이어서 $\lambda=1$이면 $P/v$가 비연결임을 보인다. 이로써 $P/v$가 연결이면 그래프도 연결이므로 정리~\ref{thm:phi}가 증명된다.

\begin{lemma}\label{lem:phi1}
설정 (S)에서 $v\in\cl(A)\cap\cl(B)$이다.
\end{lemma}
\begin{proof}
$P$가 연결이므로 모든 $x\in X$는 $v\in C$인 어떤 circuit $C$에 속하고, $C-v$는 $x$가 있는 쪽에 놓인다. 따라서 $v\in\cl(C-v)$는 그쪽의 폐포에 속한다. 양쪽 모두 공집합이 아니므로, $A$의 원소와 $B$의 원소에 각각 적용하면 $v\in\cl(A)\cap\cl(B)$이다.
\end{proof}

\begin{lemma}\label{lem:A}
설정 (S)에서 모든 $g\in A\cap\cl(B)$는 $v$와 평행하다. $B$에 대해서도 대칭적으로 성립한다.
\end{lemma}
\begin{proof}
$I$를 $B$의 기저라 하자. 보조정리~\ref{lem:phi1}에 의해 $v\in\cl(I)\setminus I$이고, $g\in\cl(I)\setminus I$이다. 기본 circuit을 $C_I(v)=D_v+v$, $C_I(g)=D_g+g$라 하면 $\varnothing\ne D_v,D_g\subseteq I$이다. $v$와 $g$를 포함하는 circuit $C$를 잡으면 $C-v\subseteq A$이다. $v\notin C_I(g)\subseteq I+g$이므로, $C$와 $C_I(g)$에서 $v$를 남기고 $g$를 강한 circuit 소거로 제거하면 $v\in C'\subseteq(C\cup C_I(g))-g$인 circuit $C'$을 얻는다. $C'-v\subseteq A$이면, $C_I(g)-g=D_g\subseteq B$이므로 $C'\subseteq C-g$가 되어 $C$가 circuit이라는 데 모순이다. 따라서 $C'-v\subseteq B$이고, $C$는 $B$와 만나지 않으므로 $C'-v\subseteq D_g$이다. $C'\subseteq I+v$이므로 $C'=C_I(v)$이고 $D_v\subseteq D_g$이다. $x\in D_v$를 고르자. $x\in D_g$는 $g$의 기본 circuit에 속하므로 $I'=I-x+g$는 독립이고 $\cl(I')=\cl(I)\ni v$이므로, $I'+v$는 유일한 circuit $C''$을 포함하며 $v\in C''$이다. $g\notin C''$이면 $C''\subseteq I+v$에서 $C''=C_I(v)\ni x$가 되어 불가능하다. 따라서 $g\in C''$이고, $C''-v\subseteq A$인 한편 $C''-\{v,g\}\subseteq I\subseteq B$이다. 그러므로 $C''=\{v,g\}$이다.
\end{proof}

\begin{proposition}\label{prop:Q}
설정 (S)에서 $\lambda=1$이다.
\end{proposition}
\begin{proof}
보조정리~\ref{lem:phi1}과 부분모듈성에 의해 $r(A)+r(B)=r(A+v)+r(B+v)\ge r(P)+r(\{v\})=r(X)+1$이므로 $\lambda\ge1$이다. (S)를 만족하면서 $\lambda\ge2$인 $(P,v,A,B)$가 있다고 하고, 그중 $|E(P)|$가 최소인 것을 고르자.

$A=\{a\}$이면 $r(A)=1$이고 $\lambda\le1+r(B)-r(X)\le1$이다. 따라서 $|A|\ge2$이고, 마찬가지로 $|B|\ge2$이다. 다음 세 가지 관찰을 쓴다. (1) $e\in A$에 대해 조건 (S)는 $P\setminus e$와 $P/e$에서도 ($A$ 대신 $A-e$로) 그대로 성립한다. minor에서 $v$를 지나는 circuit은 $v$를 지나는 $P$의 circuit에 포함되기 때문이다. (2) $e\in A$에 대해 $r_{P/e}(A-e)=r(A)-1$, $r_{P/e}(B)=r(B+e)-1$, $r_{P/e}(X-e)=r(X)-1$이므로 $\lambda(P/e)=\lambda-[e\in\cl(B)]$이다(여기서 $[\,\cdot\,]$는 명제가 참이면 $1$, 거짓이면 $0$이다). (3) $P|A$의 코루프가 아닌 $e\in A$에 대해서는 $r(A-e)=r(A)$이고 따라서 $r(X-e)=r(X)$이므로 $\lambda(P\setminus e)=\lambda$이다. 따라서 $P/e$ 또는 $P\setminus e$가 연결이고 $\lambda$가 보존되면 더 작은 반례를 얻는다.

\emph{1단계: $A$와 $B$는 독립이다.} $e\in A$라 하자. $e$가 $v$와 평행하면 $v$는 $P/e$의 루프이므로 $P/e$는 비연결이다. 그렇지 않으면 보조정리~\ref{lem:A}에 의해 $e\notin\cl(B)$이므로 (2)에 의해 $\lambda(P/e)=\lambda$이고, 최소성에 의해 $P/e$는 비연결이어야 한다. Tutte 정리에 의해 $P\setminus e$는 연결이므로, (3)과 최소성에 의해 $e$는 $P|A$의 코루프여야 한다. 따라서 $A$는 독립이고, $B$도 마찬가지이다.

\emph{2단계: 모순.} $A$가 독립이므로 $A$에서 $v$와 평행한 원소는 많아야 하나이고 $|A|\ge2$이다. $v$와 평행하지 않은 $e\in A$를 잡자. 1단계에 의해 $P/e$는 비연결이다. $W$를 $v$를 포함하는 $P/e$의 연결 성분이라 하고 $R=E(P)\setminus(W+e)$로 두면 $R\ne\varnothing$이다. $y\in R$에 대해 $v$와 $y$를 지나는 $P$의 circuit $C$를 잡는다. $e\in C$이면 $C-e$는 $v$와 $y$를 포함하는 $P/e$의 circuit이고, $e\notin\cl(C)$이면 $C$ 자체가 $P/e$의 circuit이다(둘 다 $P/e$의 circuit에 관한 표준적인 사실이다. \cite[3장]{Ox11} 참고). 두 경우 모두 $y\in W$가 되어 모순이다. 따라서 $e\in\cl(C)\setminus C$이다. $C-v\subseteq B$이면 보조정리~\ref{lem:phi1}에 의해 $e\in\cl(B)$이고, 보조정리~\ref{lem:A}에 의해 $e$는 $v$와 평행하게 되어 가정에 어긋난다. 그러므로 $C-v\subseteq A$이고, 따라서 $R\subseteq A-e$이다. $A$가 독립이므로 $r_{P/e}(R)=|R|$이고, $P/e=(P/e)|W\oplus(P/e)|R$에서 $r_{P/e}(W)=r(P/e)-|R|$이다. $e$는 루프가 아니므로 $r(W+e)=r_{P/e}(W)+1$이고 $r(P/e)=r(P)-1$이다. $W+e=E(P)\setminus R$이므로 $r(E(P)\setminus R)=r(P)-|R|$을 얻는다. 원소 하나를 삭제하면 계수는 많아야 $1$만큼 줄어들므로 $R$의 모든 원소 $x$는 $P$의 코루프이다($x$를 먼저 삭제하면, 나머지 $|R|-1$번의 삭제로는 계수가 많아야 $|R|-1$만큼 줄어든다). 이는 $P$가 연결이라는 데 모순이다.
\end{proof}

\begin{proof}[정리~\ref{thm:phi}의 증명]
그래프가 비연결이라 하고, 그 연결 성분 하나를 $A$, 나머지를 $B$로 두어 설정 (S)를 얻자. 명제~\ref{prop:Q}에 의해 $\lambda=1$이다. 보조정리~\ref{lem:phi1}에 의해 $r_{P/v}(A)=r(A)-1$, $r_{P/v}(B)=r(B)-1$이고, $r(P/v)=r(X)-1$이다. 따라서 $r_{P/v}(A)+r_{P/v}(B)=r(A)+r(B)-2=r(P/v)$이므로 $P/v=(P/v)|A\oplus(P/v)|B$는 비연결이다.
\end{proof}

\begin{remark}[Lehman 정리와의 관계]\label{rem:lehman}
정리~\ref{thm:phi}는 Lehman 정리~\cite{Le64}로부터도 유도할 수 있다(\cite[\S4.3]{Ox11} 참고). 이 정리는 다음과 같다. 연결 매트로이드 $P$에서 $v$를 지나는 서로 다른 두 circuit $C_1,C_2$마다 $(C_1\cup C_2)-\bigcap\{C: v\in C\subseteq C_1\cup C_2\}$ 꼴의 집합을 만들면, 이러한 집합들 가운데 극소인 것들이 정확히 $v$를 포함하지 않는 circuit이다.

이로부터 정리~\ref{thm:phi}를 얻는 과정은 다음과 같다. 먼저 $P/v$의 circuit은 $P$의 circuit $C$에서 $v$를 뺀 집합 $C-v$ 가운데 공집합이 아니면서 극소인 것들이다. 설정 (S)에서는 $v$를 지나는 circuit $C$마다 $C-v$가 $A$와 $B$ 중 한쪽에만 들어간다. 따라서 $A$와 $B$를 모두 만나는 $P/v$의 circuit $D$는 $v$를 포함하지 않는 $P$의 circuit 그 자체이다. 또한 $v\notin\cl(D)$이다. $v\in\cl(D)$라면 $D+v$는 $v$를 포함하는 circuit $C$를 포함하고, $D$가 circuit이므로 $C-v$는 $D$의 진부분집합이 되어 $D$의 극소성에 모순이기 때문이다.

이제 Lehman 정리에 의해 $D$를 위의 꼴로 쓰자. $C_1-v$와 $C_2-v$가 같은 쪽에 있으면 $D$도 그쪽에 들어가므로, $D$가 $A$와 $B$를 모두 만난다는 데 모순이다. 서로 다른 쪽에 있다면, (S)에 의해 $C_1\cup C_2$ 안에서 $v$를 지나는 circuit은 $C_1$과 $C_2$뿐이므로 교집합은 $\{v\}$이다. 이때 $D=(C_1\cup C_2)-\{v\}$는 $C_1-v$를 포함하므로 $v\in\cl(D)$가 되어, 앞에서 보인 사실에 모순이다.

따라서 $A$와 $B$를 모두 만나는 $P/v$의 circuit은 없으며, $P/v$는 비연결이다. 본문에서는 Lehman 정리 대신 circuit 소거, 부분모듈성, Tutte 정리만 사용하는 자기완결적인 증명을 택했으며, 이 증명은 명제~\ref{prop:Q}의 계수 등식 $\lambda=1$도 함께 보여 준다.
\end{remark}

%====================================================================
\section{균등 매트로이드}\label{sec:uniform}
균등 매트로이드에서는 정리~\ref{thm:main}의 증명에 나오는 구성을, 선택이 전혀 필요 없는 명시적 정사각형 기저로 대신할 수 있다. $[n]$의 $k$\nobreakdash-부분집합들을 colex 순서로 정렬한다. 즉 $\max(X\triangle Y)\in Y$이면 $X\colex Y$이다. $(k-1)$\nobreakdash-부분집합에도 같은 순서를 쓴다. $j$를 포함하지 않는 가장 작은 $(k-1)$\nobreakdash-부분집합은 $j\le k-1$이면 $[k]-j$이고, $j\ge k$이면 $[k-1]$이다.

\begin{lemma}[단위삼각성]\label{lem:unitri}
$S$를 $N-\rank\AM$개의 기저로 이루어진 집합이라 하고, 각 $s\in S$에 대해 $s$\nobreakdash-좌표가 $\pm1$이고 지지집합의 원소 중 $S$에 속하는 것이 모두 colex 순서로 $s$ 이하인 정사각형 $r_s$가 존재한다고 하자. 그러면 $\{r_s:s\in S\}$는 $\kerZ$의 $\Z$\nobreakdash-기저이다.
\end{lemma}
\begin{proof}
$\pi:\kerZ\to\Z^S$를 사영이라 하자. 벡터 $\pi(r_s)$들을 colex 순서로 나열하면 대각 성분이 $\pm1$인 삼각 행렬이 되므로, $L=\langle r_s\rangle$에 대해 $\pi(L)=\Z^S$이다. $\kerZ$의 계수가 $|S|$이고 $\pi$가 전사이므로 $\ker\pi$의 계수는 $0$이고, 따라서 $\ker\pi$는 자명하다. $x\in\kerZ$에 대해 $\pi(l)=\pi(x)$인 $l\in L$을 고르면 $x=l$이다. 따라서 $L=\kerZ$이고, $|S|$개의 생성원은 기저를 이룬다.
\end{proof}

\begin{proposition}\label{prop:uniform}
$1\le k\le n-1$이면 $\rank A_{U_{k,n}}=n$이고, 아래와 같이 정한 정사각형 $r_s$ ($s\in S$)들은 $\ker_\Z A_{U_{k,n}}$의 $\Z$\nobreakdash-기저를 이룬다.
\[
T=\{[k]\}\cup\{[k-1]+i:k<i\le n\}\cup\{([k]-j)+(k+1):1\le j\le k-1\}
\]
을 $n$개의 $k$\nobreakdash-부분집합으로 이루어진 집합, $S$를 그 여집합이라 하자. $s\in S$에 대해 $i=\max s$, $B=s-i$, $W_s=[i-1]\setminus B$로 둔다. $W_s$가 $k$ 이상인 원소를 포함하면 그중 가장 작은 원소를 $j$로 하고 $B'=[k-1]$로 두며, 그렇지 않으면 $j=\min W_s$로 하고 $B'=[k]-j$로 둔다. 이때 $r_s=z_{i,j;B,B'}$으로 둔다.
\end{proposition}
\begin{proof}
$M=U_{k,n}$이라 하자. $\AM\AM^{T}$의 대각 성분은 $\binom{n-1}{k-1}$이고 비대각 성분은 $\binom{n-2}{k-2}$이므로, 고윳값 $\binom{n-2}{k-1}$($\ind_n^{\perp}$ 위)과 $\binom{n-1}{k-1}+(n-1)\binom{n-2}{k-2}$가 모두 양수이고 $\rank\AM=n$이다. $T$의 세 부분은 서로소이므로 $|S|=N-n$이다.

$s\in S$라 하자. $s\ne[k]$이므로 $i>k$이고 $|W_s|=i-k\ge1$이다. 첫 번째 경우에는 $[k-1]+i\in T$이므로 $B\ne[k-1]$이고, 따라서 $B'\colex B$이며 $j\notin B'$이다. 두 번째 경우에 $B=B'=[k]-j$이면 $W_s=\{j\}\cup\{k+1,\dots,i-1\}\subseteq[k-1]$에서 $i=k+1$이 되어 $s\in T$이다. 따라서 $B\ne B'$이고, $B'$은 $j$를 포함하지 않는 가장 작은 $(k-1)$\nobreakdash-집합이므로 $B'\colex B$이다. 두 경우 모두 $j<i$, $j\notin B\cup B'$, $B'\colex B$이므로 $\max B'\le\max B<i$이고 $i\notin B'$이다. 모든 $k$\nobreakdash-집합이 기저이므로 $r_s=e_{B+i}-e_{B+j}-e_{B'+i}+e_{B'+j}$는 정사각형이다. $B'+i\colex B+i$이고 $B+j,B'+j\colex B+i$(최댓값이 $i$보다 작다)이므로, 지지집합에서 colex 순서로 가장 큰 원소는 계수 $+1$인 $s=B+i$이다. 보조정리~\ref{lem:unitri}로 증명이 끝난다.
\end{proof}

\begin{example}\label{ex:u24}
$U_{2,4}$에서는 $T=\{12,13,14,23\}$, $S=\{24,34\}$이고, 그림~\ref{fig:u24}의 기호로 $r_{24}=e_{24}-e_{23}-e_{14}+e_{13}=z_1$, $r_{34}=e_{34}-e_{23}-e_{14}+e_{12}=z_2$이다.
\end{example}


%====================================================================
\section{귀결}\label{sec:apps}
\subsection{토릭 아이디얼}
$K$를 체, $K[y]=K[y_I:I\in\bs]$, $p=\prod_{I\in\bs}y_I$, $\mathcal Y=K[y][p^{-1}]=K[y^{\pm1}]$이라 하자. $z\in\Z^{\bs}$에 대해 $f_z=y^{z^+}-y^{z^-}$로 둔다. $z$가 정사각형이면 $f_z$는 대칭 교환 이항식이다(보조정리~\ref{lem:sym}).

\begin{corollary}\label{cor:toric}
$M$을 계수가 $1$ 이상인 매트로이드, $c=N-\rank\AM=\operatorname{ht}I_M$이라 하자. $c$개의 정사각형으로 이루어진 집합 $\mathcal R$에 대해 다음은 동치이다.
\begin{enumerate}[label=(\alph*),itemsep=1pt]
\item $\mathcal R$은 $\kerZ$의 $\Z$\nobreakdash-기저이다.
\item $I_M\,\mathcal Y=(f_z:z\in\mathcal R)\,\mathcal Y$이다.
\item $I_M=(f_z:z\in\mathcal R):p^{\infty}$이다.
\end{enumerate}
특히 정리~\ref{thm:main}에 의해 이러한 $\mathcal R$이 존재하므로, $I_M$은 $c$개의 대칭 교환 이항식이 생성하는 아이디얼의 $p$\nobreakdash-포화이며, 생성원이 $c$개보다 적은 아이디얼의 $p$\nobreakdash-포화는 $I_M$이 될 수 없다.
\end{corollary}
\begin{proof}
$\mathcal Y$에서 $f_z=y^{z^-}(y^z-1)$이므로 $(f_z:z\in\mathcal R)\mathcal Y=(y^z-1:z\in\mathcal R)\mathcal Y$이다. $y^{v+w}-1=y^v(y^w-1)+(y^v-1)$이고 $y^{-v}-1=-y^{-v}(y^v-1)$이므로, $L=\langle\mathcal R\rangle$이라 하면 이 아이디얼은 $\mathfrak a_L=(y^u-1:u\in L)\mathcal Y$와 같고, 마찬가지로 $I_M\mathcal Y=\mathfrak a_{\kerZ}$이다. 사상 $y^w\mapsto[w]$는 군대수로의 동형사상 $\mathcal Y/\mathfrak a_L\cong K[\Z^{\bs}/L]$을 유도하며, 이때 $y^u-1$은 $[u]-[0]$으로 간다. 서로 다른 군 원소는 일차독립이므로 $y^u-1\in\mathfrak a_L$일 필요충분조건은 $u\in L$이다. 따라서 (b)는 $L=\kerZ$과 동치이고, $|\mathcal R|=c=\rank\kerZ$이므로 이는 (a)와 동치이다.

임의의 아이디얼 $J\subseteq K[y]$에 대해 $J\mathcal Y\cap K[y]=J:p^{\infty}$이다(이는 표준적인 사실이다. \cite{ES96}을 보라). 토릭 아이디얼은 소 아이디얼이고 단항식을 포함하지 않으므로 $I_M:p^\infty=I_M$이다. 따라서 (b)가 성립하면 $(f_z:z\in\mathcal R):p^\infty=I_M\mathcal Y\cap K[y]=I_M$이므로 (c)가 성립한다. 역으로 (c)가 성립하면 $(f_z)\mathcal Y=((f_z):p^\infty)\mathcal Y=I_M\mathcal Y$이므로 (b)가 성립한다.

끝으로 $J\subseteq K[y]$를 $J:p^\infty=I_M$인 임의의 아이디얼이라 하자. 위와 같이 $J\mathcal Y=I_M\mathcal Y$이다. 국소화는 곱셈적 집합과 만나지 않는 소 아이디얼의 높이를 보존하므로 $J\mathcal Y$는 높이가 $c$인 소 아이디얼이고, 따라서 자기 자신 위의 유일한 극소 소 아이디얼이다. Krull 높이 정리에 의해 생성원이 $c$개보다 적은 아이디얼 위의 극소 소 아이디얼은 높이가 $c$보다 작으므로, $J\mathcal Y$는 $c$개보다 적은 원소로 생성될 수 없다. $J$의 생성원들은 $J\mathcal Y$를 생성하므로 $J$에도 적어도 $c$개의 생성원이 필요하다.
\end{proof}

\begin{remark}
$J_M$을 모든 대칭 교환 이항식이 생성하는 아이디얼이라 하고 $\mathfrak m=(y_I:I\in\bs)$라 하자. \LM{} 정리(정리~\ref{thm:sat})에 따르면 $J_M:\mathfrak m^\infty=I_M$이다. $p\in\mathfrak m$이므로 $J_M:\mathfrak m^\infty\subseteq J_M:p^\infty$이고, $J_M\subseteq I_M$이며 $I_M$은 $p$를 포함하지 않는 소 아이디얼이므로 $J_M:p^\infty\subseteq I_M$이다. 따라서 $J_M:p^\infty=I_M$이다. 따름정리~\ref{cor:toric}은 정사각형 기저에 해당하는 $c$개의 대칭 교환 이항식만으로 이미 이 등식이 성립하며, 생성원이 $c$개보다 적은 아이디얼로는 성립할 수 없음을 보여 준다.
\end{remark}

\begin{remark}[Lattice basis와 Markov basis]\label{rem:markov}
대수 통계학에서 $\kerZ$의 유한 부분집합은, 대응하는 이항식들이 $I_M$을 생성하면 \emph{Markov basis}라 하고, $\kerZ$의 $\Z$\nobreakdash-기저이면 \emph{lattice basis}라 한다~\cite{DS98}. Markov basis는 항상 $\kerZ$을 군으로서 생성하지만, lattice basis는 일반적으로 Markov basis가 아니며~\cite{DS98}, $\kerZ$을 생성하는 집합이 lattice basis를 포함한다는 보장도 없다. Lattice basis의 이항식들이 생성하는 아이디얼을 \emph{lattice basis ideal}이라 한다. 이 아이디얼의 $p$\nobreakdash-포화는 격자 아이디얼, 즉 여기서는 $I_M$이며~\cite{ES96}, lattice basis ideal의 준소 분해는 Ho{\lmfont ş}ten과 Shapiro가 연구하였다~\cite{HS00}. 보조정리~\ref{lem:sym}에 의해 White 추측은 정확히 정사각형들이 Markov basis를 이룬다는 주장이며, 이는 일반적으로 성립하지 않는다~\cite{Lar26}. 반면 정리~\ref{thm:main}은 계수가 $1$ 이상이면 정사각형들이 항상 lattice basis를 포함함을 보여 주고, 따름정리~\ref{cor:toric}은 이에 대응하는 lattice basis ideal의 $p$\nobreakdash-포화가 $I_M$임을 보여 준다.
\end{remark}

\subsection{가역 열 스케일링}\label{sec:criterion}
Plücker 좌표에 대한 열 토러스의 작용은 고전적이며~\cite{GGMS87,DW91}, 행 스케일링과 열 스케일링을 매트로이드에 연관된 토릭 아이디얼과 함께 다루는 틀로는 slack realization space가 있다~\cite{BW19}. 아래 판정법은 이 틀에 속하는 응용이며, 정리~\ref{thm:main}의 정수성이 실제로 쓰이는 곳이기도 하다. 계수가 $k$인 $k\times n$ 실수 행렬 $Q,Q'$에 대해, 어떤 $G\in GL(k,\R)$과 가역 대각 행렬 $D$로 $Q'=GQD$가 되는지, 그리고 대각 성분이 양수인 대각 행렬 $D$로 그렇게 되는지를 묻는다. $k$\nobreakdash-부분집합 $I$에 대해 $Q_I$를 $I$에 속하는 열들로 이루어진 $Q$의 $k\times k$ 부분행렬이라 하고, $p_I(Q)=\det Q_I$, $\bs(Q)=\{I:p_I(Q)\ne0\}$으로 둔다. $D=\operatorname{diag}(\mu_1,\dots,\mu_n)$에 대해
\begin{equation}\label{eq:plucker}
p_I(GQD)=\det(G)\Big(\prod_{i\in I}\mu_i\Big)p_I(Q)
\end{equation}
이다. $\bs(Q)=\bs(Q')=\bs$일 때 $I\in\bs$에 대해 $\epsilon_I=\sgn\big(p_I(Q')/p_I(Q)\big)$, $\eta_I=\log\big|p_I(Q')/p_I(Q)\big|$로 둔다.

\begin{theorem}\label{thm:criterion}
$\bs(Q)=\bs(Q')=\bs$이고, $M$을 기저 집합이 $\bs$인 매트로이드, $\mathcal R$을 $M$의 정사각형 기저라 하자.
\begin{enumerate}[label=(\alph*),itemsep=1pt]
\item 어떤 $G\in GL(k,\R)$과 가역 대각 행렬 $D$에 대해 $Q'=GQD$일 필요충분조건은, 모든 $z\in\mathcal R$에 대해
\[
\prod_{I\in\bs}\epsilon_I^{\,z_I}=1\qquad\text{이고}\qquad z^{T}\eta=0
\]
인 것이다.
\item 어떤 $G\in GL(k,\R)$과 양의 대각 행렬 $D$에 대해 $Q'=GQD$일 필요충분조건은, $\epsilon_I$가 $\bs$ 위에서 상수이고 모든 $z\in\mathcal R$에 대해 $z^{T}\eta=0$인 것이다.
\end{enumerate}
\end{theorem}
\begin{proof}
\emph{필요성.} $D=\operatorname{diag}(\mu_1,\dots,\mu_n)$에 대해 $Q'=GQD$라 하자. \eqref{eq:plucker}에 의해 $\epsilon_I=\sgn\det G\prod_{i\in I}\sgn\mu_i$이고 $\eta_I=\log|\det G|+\sum_{i\in I}\log|\mu_i|$이다. $z\in\kerZ$이면 $z^{T}\ind_{\bs}=\tfrac1k\ind_n^{T}\AM z=0$이므로
\begin{gather*}
\prod_{I}\epsilon_I^{\,z_I}=(\sgn\det G)^{z^T\ind_{\bs}}\prod_i(\sgn\mu_i)^{(\AM z)_i}=1,\\
z^T\eta=\log|\det G|\, z^T\ind_{\bs}+\sum_i\log|\mu_i|\,(\AM z)_i=0
\end{gather*}
이다. $D$가 양의 대각 행렬이면 $\epsilon_I=\sgn\det G$는 상수이다.

\emph{절댓값.} 모든 $z\in\mathcal R$에 대해 $z^T\eta=0$이면, $\mathcal R$이 $\kerR$을 생성하므로 $\eta$는 $\kerR$에 직교하고, 따라서 어떤 $x\in\R^n$에 대해 $\eta=\AM^{T}x$이다. 즉 모든 $I\in\bs$에 대해 $|p_I(Q')|=|p_I(Q)|\prod_{i\in I}e^{x_i}$이다.

\emph{(b)의 충분성.} $D=\operatorname{diag}(e^{x_i})$로 두면, 공통 부호를 $\epsilon$이라 할 때 Plücker 벡터로서 $p(Q')=\epsilon\,p(QD)$이다. Plücker 매장의 단사성에 의해 $Q'$과 $QD$는 같은 행공간을 가지므로 어떤 $G\in GL(k,\R)$에 대해 $Q'=GQD$이고, $\det G=\epsilon$이다.

\emph{(a)의 충분성.} $\epsilon_I=(-1)^{s_I}$인 $s_I\in\Z/2$를 잡자. $\mathcal R$은 $\kerZ$의 $\Z$\nobreakdash-기저이므로, 부호 조건에서 모든 $u\in\kerZ$에 대해 $\Z/2$에서 $\sum_I u_Is_I=0$이다. 따라서 $\varphi(\AM u)=\sum_I u_Is_I$는 $\Lambda=\AM(\Z^{\bs})$ 위의 준동형사상 $\varphi:\Lambda\to\Z/2$로 잘 정의된다. $V^0\subseteq\Z^{\bs}$를 좌표합이 $0$인 벡터들의 부분격자, $\Lambda^0=\AM(V^0)$이라 하자. 명제~\ref{prop:dsum}의 증명 3단계에서처럼, $\Lambda^0$은 벡터 $e_a-e_b$들로 이루어진 $\Z$\nobreakdash-기저를 가진다. 여기서 $\{a,b\}$는 $F$의 변 전체에 걸치고, $F$는 어떤 $W$에 대해 $W+a,W+b\in\bs$일 때 $a\sim b$로 정한 $[n]$ 위의 그래프의 spanning forest이다. $F$의 모든 변 $\{a,b\}$에 대해 $\sigma_a-\sigma_b=\varphi(e_a-e_b)$인 $\sigma\in(\Z/2)^n$을 고르자(각 나무의 한 꼭짓점에서 $\sigma$를 정하고 나무를 따라 전파한다). 그러면 모든 $t\in\Lambda^0$에 대해 $\varphi(t)=\sum_i\sigma_it_i$이다. $J_0\in\bs$를 고정하고 $\tau=\varphi(\ind_{J_0})-\sum_{i\in J_0}\sigma_i$로 두자. $w=\AM u\in\Lambda$이고 $m=\sum_Iu_I$이면 $w-m\ind_{J_0}=\AM(u-me_{J_0})\in\Lambda^0$이므로 $\varphi(w)=\sum_i\sigma_iw_i+m\tau$이다. $w=\ind_I$로 두면 $s_I=\tau+\sum_{i\in I}\sigma_i$, 즉 $\epsilon_I=(-1)^{\tau}\prod_{i\in I}(-1)^{\sigma_i}$를 얻는다. 이제 $D=\operatorname{diag}\big((-1)^{\sigma_i}e^{x_i}\big)$로 두면 모든 $I\in\bs$에 대해
\[
p_I(QD)=p_I(Q)\prod_{i\in I}(-1)^{\sigma_i}e^{x_i}=(-1)^{\tau}\epsilon_I\,|p_I(Q')|\,\frac{p_I(Q)}{|p_I(Q)|}=(-1)^{\tau}p_I(Q')
\]
이고, $I\notin\bs$이면 양변이 모두 $0$이다. 앞에서와 같이 $\det G=(-1)^{\tau}$인 어떤 $G\in GL(k,\R)$에 대해 $Q'=GQD$이다.
\end{proof}

따라서 (b)에서 부호가 상수라는 조건을 제외하면, $N$개의 Plücker 비 전체에 대한 조건이 정사각형 기저를 이루는 $N-\rank\AM$개의 정사각형마다 하나씩 주어지는 조건으로 줄어든다. $M$이 연결이면 $D$는 (a)에서는 $0$이 아닌 상수배를 제외하면, (b)에서는 양의 상수배를 제외하면 유일하다. 절댓값은 $\rank\AM=n$이므로(보조정리~\ref{lem:lat}) 공통 인자를 제외하면 정해진다. 또한 위 그래프에서 임의의 두 원소 $a,b$에 대해 $a\sim b$이므로(보조정리~\ref{lem:lat}의 증명) $\sigma_a-\sigma_b=s_{W+a}-s_{W+b}$가 강제되고, 따라서 부호도 공통 부호를 제외하면 정해진다.

\begin{remark}\label{rem:index}
(b)는 $\mathcal R$이 $\kerR$을 생성한다는 사실만 쓰지만, (a)는 $\mathcal R$이 $\kerZ$을 생성한다는 사실을 쓴다. 부호 조건은 $\{\pm1\}\cong\Z/2$에 값을 갖기 때문이다. 정사각형으로 이루어진 $\kerR$의 $\R$\nobreakdash-기저는 진부분격자를 생성할 수 있다. 예를 들어 $U_{2,5}$에서 다섯 정사각형 $z_{1,2;\{3\},\{4\}}$, $z_{1,3;\{2\},\{5\}}$, $z_{1,4;\{3\},\{5\}}$, $z_{1,5;\{2\},\{4\}}$, $z_{2,3;\{4\},\{5\}}$는 일차독립이고 $\rank\ker_\Z A_{U_{2,5}}=10-5=5$이지만, 이들은 지수가 $2$인 부분격자를 생성한다.
\end{remark}

%====================================================================
\section{추가 논의와 열린 문제}\label{sec:caveats}
\begin{theorem}[\LM{}~\cite{LM14}]\label{thm:sat}
계수가 $1$ 이상인 모든 매트로이드 $M$과 모든 체 $K$에 대해 $I_M=J_M:\mathfrak m^\infty$이다. 여기서 $J_M$은 대칭 교환 이항식들이 생성하는 아이디얼이고 $\mathfrak m=(y_I:I\in\bs)$이다.
\end{theorem}

\begin{proposition}\label{prop:lattice}
$\rank M\ge1$이면 $\sq(M)=\kerZ$이다. 이는 정리~\ref{thm:main}에서 따라 나오며, 정리~\ref{thm:sat}에서도 직접 따라 나온다.
\end{proposition}
\begin{proof}[정리~\ref{thm:sat}을 이용한 증명]
$L=\sq(M)$, $u\in\kerZ$이라 하자. 정리~\ref{thm:sat}과 보조정리~\ref{lem:sym}에 의해 어떤 기저 $I_0$과 어떤 $m$에 대해 $y_{I_0}^{m}(y^{u^+}-y^{u^-})\in J_M\subseteq(f_z:z\in L)$이다. 따름정리~\ref{cor:toric}의 증명에서처럼 $\mathcal Y$와 $K[\Z^{\bs}/L]$로 넘어가면 $[me_{I_0}]([u^+]-[u^-])=0$을 얻는다. $[me_{I_0}]$은 군 원소이므로 가역이고, 따라서 $[u^+]=[u^-]$, 즉 $u\in L$이다.
\end{proof}

\begin{remark}[White 추측]\label{rem:white}
정리~\ref{thm:main}과 따름정리~\ref{cor:toric}은 격자 $\kerZ$, 즉 모든 변수를 가역으로 만든 뒤의 아이디얼 $I_M$에 관한 것이다. White 추측은 $K[y]$ 자체에서 $J_M=I_M$인지를 묻는다. 이는 격자가 아니라 아이디얼에 관한 진술이며, 일반적으로는 거짓이다~\cite{Lar26}.
\end{remark}

\medskip\noindent\textbf{열린 문제.}
\begin{enumerate}[label=(\arabic*),itemsep=2pt]
\item 명제~\ref{prop:uniform}를 일반화하는 정준적인 정사각형 기저, 예를 들어 기저 교환 그래프나 $[n]$ 위의 전순서로 결정되는 정사각형 기저가 있는가?
\item 어떤 매트로이드에서 $(f_z:z\in\mathcal R):\mathfrak m^\infty=I_M$인 정사각형 기저 $\mathcal R$이 존재하는가? $\mathfrak m$\nobreakdash-포화는 $p$\nobreakdash-포화보다 세밀하므로, 이는 따름정리~\ref{cor:toric}을 정리~\ref{thm:sat} 쪽으로 강화하게 된다.
\item 독립성 오라클이 주어졌을 때 $n$과 $N$에 대한 다항 시간 안에 정사각형 기저를 출력하는 알고리즘이 있는가?
\end{enumerate}

\begin{thebibliography}{99}
\bibitem{W80} N. White, \emph{A unique exchange property for bases}, Linear Algebra Appl. 31 (1980), 81--91.
\bibitem{Ka10} K. Kashiwabara, \emph{The toric ideal of a matroid of rank 3 is generated by quadrics}, Electron. J. Combin. 17 (2010), \#R28.
\bibitem{Bl08} J. Blasiak, \emph{The toric ideal of a graphic matroid is generated by quadrics}, Combinatorica 28 (2008), 283--297.
\bibitem{Sc11} J. Schweig, \emph{Toric ideals of lattice path matroids and polymatroids}, J. Pure Appl. Algebra 215 (2011), 2660--2665.
\bibitem{Bo13} J. Bonin, \emph{Basis-exchange properties of sparse paving matroids}, Adv. in Appl. Math. 50 (2013), 6--15.
\bibitem{YY25} Y.-C. Yu and C. H. Yuen, \emph{White's conjecture for paving matroids}, arXiv:2510.04163 (2025).
\bibitem{LM14} M. Lasoń and M. Micha\lstroke{}ek, \emph{On the toric ideal of a matroid}, Adv. Math. 259 (2014), 1--12.
\bibitem{La21} M. Lasoń, \emph{On the toric ideals of matroids of a fixed rank}, Selecta Math. (N.S.) 27 (2021), Paper No.~18.
\bibitem{HMW25} K. Han, M. Micha\lstroke{}ek, and J. Weigert, \emph{White's conjecture for matroids and inner projections}, arXiv:2501.17738 (2025).
\bibitem{HHKS20} K. Hayase, T. Hibi, K. Katsuno, and K. Shibata, \emph{Gröbner bases of toric ideals associated with matroids}, Acta Math. Vietnam. 47 (2022), 775--779.
\bibitem{BMS24} K. Bérczi, B. Mátravölgyi, and T. Schwarcz, \emph{Reconfiguration of basis pairs in regular matroids}, J. Combin. Theory Ser. B 177 (2026), 105--142.
\bibitem{BS24} K. Bérczi and T. Schwarcz, \emph{Exchange distance of basis pairs in split matroids}, SIAM J. Discrete Math. 38 (2024), no. 1, 132--147.
\bibitem{Lar26} M. Larson, \emph{Counterexamples to two conjectures about matroids}, arXiv:2607.02208 (2026).
\bibitem{BCLLM26} S. Backman, N. Cheung, M. Lasoń, G. Liu, and M. Micha\lstroke{}ek, \emph{The Gröbner version of White's conjecture is false}, arXiv:2606.13960 (2026).
\bibitem{DFMR26} J. A. De Loera, L. Ferroni, S. Morales, and J. Rambau, \emph{There are matroid toric ideals without quadratic Gröbner bases}, arXiv:2606.11014 (2026).
\bibitem{DS98} P. Diaconis and B. Sturmfels, \emph{Algebraic algorithms for sampling from conditional distributions}, Ann. Statist. 26 (1998), 363--397.
\bibitem{W77} N. White, \emph{The basis monomial ring of a matroid}, Adv. Math. 24 (1977), 292--297.
\bibitem{St96} B. Sturmfels, \emph{Gröbner Bases and Convex Polytopes}, University Lecture Series 8, Amer. Math. Soc., 1996.
\bibitem{GR15} I. García-Marco and J. L. Ramírez Alfonsín, \emph{Matroid toric ideals: complete intersection, minors and minimal systems of generators}, SIAM J. Discrete Math. 29 (2015), no. 4, 2267--2276.
\bibitem{Ox11} J. Oxley, \emph{Matroid Theory}, 2nd ed., Oxford University Press, 2011.
\bibitem{Br69} R. A. Brualdi, \emph{Comments on bases in dependence structures}, Bull. Austral. Math. Soc. 1 (1969), 161--167.
\bibitem{Le64} A. Lehman, \emph{A solution of the Shannon switching game}, J. Soc. Indust. Appl. Math. 12 (1964), 687--725.
\bibitem{ES96} D. Eisenbud and B. Sturmfels, \emph{Binomial ideals}, Duke Math. J. 84 (1996), 1--45.
\bibitem{HS00} S. Ho{\lmfont ş}ten and J. Shapiro, \emph{Primary decomposition of lattice basis ideals}, J. Symbolic Comput. 29 (2000), 625--639.
\bibitem{GGMS87} I. M. Gelfand, R. M. Goresky, R. D. MacPherson, and V. V. Serganova, \emph{Combinatorial geometries, convex polyhedra, and Schubert cells}, Adv. Math. 63 (1987), 301--316.
\bibitem{DW91} A. Dress and W. Wenzel, \emph{Grassmann--Pl\"ucker relations and matroids with coefficients}, Adv. Math. 86 (1991), 68--110.
\bibitem{BW19} M. Brandt and A. Wiebe, \emph{The slack realization space of a matroid}, Algebr. Comb. 2 (2019), no. 4, 663--681.
\end{thebibliography}

\par\vfill
\noindent\rule{\linewidth}{0.4pt}\\[0.2em]
{\small\noindent\textbf{라이선스.} 저작권 \copyright{} 2026 김영수. 이 저작물은 크리에이티브 커먼즈 저작자표시 4.0 국제 라이선스(CC BY 4.0, \url{https://creativecommons.org/licenses/by/4.0/deed.ko})에 따라 이용할 수 있다.\par}
\end{document}
